\clearpage
\begin{examplebox}
\textbf{\textcolor{CExample}{例 1｜基础：由中点求斜率、弦方程}}
椭圆 $\frac{x^2}{16}+\frac{y^2}{12}=1$ 的一条弦中点为 $M(2,1)$，求斜率和直线方程。
\begin{center}\begin{tikzpicture}[x=1.1cm,y=1.1cm,>=stealth,font=\small]
\node[font=\small\bfseries] at (0,4.8) {例 1：中点给出斜率 $-3/2$};

\draw[->,gray] (-4.4,0)--(4.5,0) node[right] {$x$};
\draw[->,gray] (0,-3.8)--(0,4) node[above] {$y$};
\draw[CStatement,thick] (0,0) ellipse (4 and 3.464101615);
\draw[teal,very thick] (.367006838,3.449489743)--(3.632993162,-1.449489743);
\fill[orange!80!black] (2,1) circle(2pt) node[right] {$M(2,1)$};
\node at (-2,-2) {$3x+2y-8=0$};
\node at (0,-4.4) {$x^2/16+y^2/12=1$};

\end{tikzpicture}
\end{center}
\par\textbf{识别信号：}已知中点，直接点差法。先核验 $4/16+1/12=1/3<1$，确有真实弦。
\[
k=-\frac{12\times2}{16\times1}=-\frac32,\qquad
 y-1=-\frac32(x-2)\quad\Longrightarrow\quad\boxed{3x+2y-8=0}.
\]
\textbf{更快的无除法法：}
$\frac2{16}(x-2)+\frac1{12}(y-1)=0$，两边乘 $24$，直接得 $3x+2y-8=0$。
\par\textbf{答案：}$k=-\frac32$；$3x+2y-8=0$。
\medskip
\textbf{\textcolor{CExample}{例 2｜逆用：已知弦的直线，求中点}}
椭圆 $\frac{x^2}{9}+\frac{y^2}{4}=1$ 与直线 $2x-3y+1=0$ 所截弦中点为 $M(x_0,y_0)$，求 $M$。
\par\textbf{最低运算量：}已知弦斜率 $k=2/3$。直接写无除法方向关系
\[
 \frac{x_0}{9}+\frac{(2/3)y_0}{4}=0
 \ \Longrightarrow\ 2x_0+3y_0=0.
\]
再加“中点在线上”：
\[
\begin{cases}2x_0+3y_0=0,\\2x_0-3y_0+1=0,\end{cases}
\quad\Longrightarrow\quad\boxed{M(-\frac14,\frac16)}.
\]
核验 $\frac{1}{16\cdot9}+\frac{1}{36\cdot4}=\frac1{144}+\frac1{144}=\frac1{72}<1$，弦真实存在。
\medskip
\textbf{\textcolor{CExample}{例 3｜变式：已知斜率和中点横坐标}}
椭圆 $\frac{x^2}{12}+\frac{y^2}{8}=1$ 的一条弦斜率为 $2$，其中点横坐标为 $1$，求中点与弦方程。
\par 设 $M(1,y_0)$，无需先除 $y_0$，由方向关系
\[
 \frac1{12}+\frac{2y_0}{8}=0
 \ \Longrightarrow\ y_0=-\frac13.
\]
核验 $1/12+(1/9)/8=7/72<1$；由点斜式
\[
 y+\frac13=2(x-1)\ \Longrightarrow\ \boxed{6x-3y-7=0}.
\]
\textbf{答案：}$M(1,-\frac13)$；$6x-3y-7=0$。
\medskip
\textbf{\textcolor{CExample}{例 4｜特殊情形：轴上中点与中心中点}}
在椭圆 $\frac{x^2}{16}+\frac{y^2}{9}=1$ 中：
\begin{itemize}[leftmargin=*,itemsep=0.3em]
\item 若 $M(0,1)$，弦为 $y=1$，\textbf{水平}，$k=0$。
\item 若 $M(2,0)$，弦为 $x=2$，\textbf{竖直}，不存在斜率。
\item 若 $M(0,0)$，$y=0$、$x=0$、$y=x$ 等所有过中心的直线均可作中心弦；\textbf{中点不决定唯一斜率}。
\end{itemize}
\medskip
\textbf{\textcolor{CExample}{例 5｜高价值变式：方程可解，弦未必存在}}
椭圆 $\frac{x^2}{25}+\frac{y^2}{9}=1$ 上斜率为 $k=\frac34$ 的弦，其中点还满足 $x_0+2y_0=c$。问 $c=1$ 时能否存在？
\par 由中点弦方向关系：
\[
 \frac{x_0}{25}+\frac{(3/4)y_0}{9}=0
 \ \Longrightarrow\ 12x_0+25y_0=0.
\]
联立 $x_0+2y_0=c$ 得 $M=(25c,-12c)$。中点必须严格在椭圆内部：
\[
 \frac{(25c)^2}{25}+\frac{(-12c)^2}{9}=41c^2<1
 \ \Longleftrightarrow\ |c|<\frac1{\sqrt{41}}.
\]
\textbf{因此 $c=1$ 不存在真实弦}。例如 $c=\frac1{10}$ 时，$M(\frac52,-\frac65)$ 则在内部，真实弦存在。
\par\textbf{迁移信号：}点差法是\textbf{必要关系}；只算出中点还不够，必须检查内部条件。
\end{examplebox}
